\documentclass[10pt,a4paper]{article}
\usepackage[margin=0.9in]{geometry}
\usepackage[T1]{fontenc}
\usepackage[utf8]{inputenc}
\usepackage{lmodern}
\usepackage{hyperref}
\usepackage{enumitem}
\usepackage{amsmath}

\hypersetup{
  colorlinks=true,
  linkcolor=blue,
  urlcolor=blue
}

\setlength{\parindent}{0pt}
\setlength{\parskip}{4pt}
\setlist[itemize]{noitemsep, topsep=2pt}
\setlist[enumerate]{noitemsep, topsep=2pt}

\title{Research Notes: Unlearning in VLMs Pipeline}
\author{Aaditya Bhatnagar}
\date{}

\begin{document}
\maketitle

\section{Introduction}
These notes will cover everything I have studied or researched while building this project.

\section{Entry Points}
There are two entry points to start building and working with this pipeline:
\begin{itemize}
  \item \texttt{scripts/run\_pipeline.py}: 
  Running this runs one complete experiment request end-to-end, but make sure the hyperparameter config either manually through the command or through \texttt{config/parameters.yaml} is configured properly.
  \item \texttt{scripts/run\_study.py}: 
  This is an optional outer wrapper that is used to run multi-seed experiments. It launches the pipeline as a \texttt{(request, seed)} pair.
\end{itemize}

For a normal single experiment, the main control flow is:
\[
\texttt{run\_pipeline.py}
\rightarrow
\texttt{prepare\_data.py}
\rightarrow
\texttt{train\_vlm.py}
\rightarrow
\texttt{run\_unlearning.py}
\rightarrow
\texttt{evaluate\_attacks.py}
\]

\section{Stage 1: Pipeline Orchestrator}
\texttt{scripts/run\_pipeline.py} is the main orchestrating script for one full run of the pipeline. 

\texttt{scripts/run\_pipeline.py} has many cli arguments that can be passed to it. The following is the list of arguments/hyperparameters that are included for this orchestrating script. 

\begin{itemize}
  \item \texttt{---config}: Default is \texttt{config/parameters.yaml},
  \item \texttt{---data-dir},
  \item \texttt{---dataset},
  \item \texttt{---request},
  \item \texttt{---split-path},
  \item \texttt{---model-name},
  \item \texttt{---output-root},
  \item \texttt{---forget-classes},
  \item \texttt{---forget-fraction},
  \item \texttt{---batch-size},
  \item \texttt{---num-workers},
  \item \texttt{---ft-epochs},
  \item \texttt{---ft-max-steps},
  \item \texttt{---ul-steps},
  \item \texttt{---methods},
  \item \texttt{---seed},
  \item \texttt{---device},
  \item \texttt{---show-config}: This prints the fully resolved settings while not running the pipeline,
  \item \texttt{---resume}: Used to resume the pipeline execution if interrupted mid-way, checks the logs, state, and i/p o/p fingerprints,
  \item \texttt{---force-stage}: Rerun a particular stage even under ---resume,
\end{itemize}

\subsection{How stage execution is guarded}
Each stage is executed through \texttt{src/unml/pipeline.py::run\_pipeline\_stage}.

That helper:
\begin{itemize}
  \item fingerprints declared input files with SHA256,
  \item records the Git commit and exact command,
  \item fingerprints declared output files after success,
  \item stores a receipt under:
  \[
  \texttt{<output\_root>/.pipeline/stages/<stage>.json}
  \]
  \item skips the stage under \texttt{--resume} only if contract and outputs still match.
\end{itemize}

\subsection{Actual stage order}
\begin{enumerate}
  \item \texttt{prepare}
  \item \texttt{finetune}
  \item \texttt{retrain\_oracle} if enabled
  \item \texttt{unlearn:<method>} for each configured method
  \item \texttt{evaluate}
\end{enumerate}

\section{Stage 2: Data Preparation}
The command dispatched is \texttt{scripts/prepare\_data.py}.

\subsection{Internal flow}
\begin{enumerate}
  \item Load runtime config.
  \item Resolve dataset, request, split path, forget classes, seed, and request metadata.
  \item Call \texttt{src/unml/data.py::download\_and\_prepare\_splits}.
  \item Load or download the CIFAR train/test datasets.
  \item Validate the hierarchy request if a named CIFAR-100 request is used.
  \item Build split indices using \texttt{make\_splits}.
  \item Save the split JSON.
\end{enumerate}

\subsection{Important split semantics}
The split JSON stores more than just forget and retain partitions. It also stores hierarchy-aware groups used later by H-TGSD and evaluation:
\begin{itemize}
  \item \texttt{forget\_indices}
  \item \texttt{retain\_train\_indices}
  \item \texttt{retain\_val\_indices}
  \item \texttt{finetune\_train\_indices}
  \item \texttt{sibling\_retain\_train\_indices}
  \item \texttt{unrelated\_retain\_train\_indices}
  \item \texttt{test\_forget\_indices}
  \item \texttt{test\_sibling\_indices}
  \item \texttt{test\_unrelated\_indices}
  \item \texttt{test\_retain\_indices}
  \item \texttt{test\_all\_indices}
  \item \texttt{target\_classes}, \texttt{sibling\_classes}, and \texttt{unrelated\_classes}
\end{itemize}

\subsection{Key code dependency}
H-TGSD only works because this stage writes hierarchy-aware metadata into the split JSON. If the request type or sibling metadata is missing, \texttt{run\_unlearning.py} will later fail validation for H-TGSD.

\section{Stage 3: Fine-Tuning}
The command dispatched is \texttt{scripts/train\_vlm.py}.

\subsection{Internal flow}
\begin{enumerate}
  \item Resolve data, split, model, adapter, optimizer, and device settings.
  \item Build a \texttt{FineTuneConfig}.
  \item Call \texttt{src/unml/train.py::run\_finetuning}.
\end{enumerate}

\subsection{What \texttt{run\_finetuning} does}
\begin{enumerate}
  \item Validate mode: either standard \texttt{finetune} or \texttt{retrain\_oracle}.
  \item Load split metadata and class names.
  \item Build CLIP image processor and tokenizer.
  \item Build class prompts with \texttt{build\_text\_inputs}.
  \item Build dataloaders with \texttt{build\_loaders}.
  \item Build or restore \texttt{LightweightVLM} from \texttt{src/unml/model.py}.
  \item Save the initialization checkpoint immediately.
  \item Create AdamW optimizer and cosine scheduler.
  \item Train on either:
  \begin{itemize}
    \item \texttt{finetune\_train} for normal fine-tuning, or
    \item \texttt{retain\_train} for retraining oracle.
  \end{itemize}
  \item Evaluate after each epoch using \texttt{retain\_val}, \texttt{test\_all}, \texttt{test\_retain}, and forget-train accuracy.
  \item Save the best checkpoint by \texttt{retain\_val\_acc}.
\end{enumerate}

\subsection{Model behavior}
\texttt{LightweightVLM} loads a frozen CLIP backbone and only exposes trainable adapter-style parameters:
\begin{itemize}
  \item \texttt{post\_projection} mode: low-rank residual adapters after image and text projection.
  \item \texttt{vision\_lora} mode: LoRA inserted into CLIP vision-attention \texttt{q\_proj} and \texttt{v\_proj}.
\end{itemize}

The text encoder is frozen. Depending on adapter type, text features may be cacheable during training. This matters later because H-TGSD requires frozen-text, cacheable text features.

\subsection{Fine-tuning artifacts}
Typical outputs are:
\begin{itemize}
  \item \texttt{finetune/checkpoints/base\_init.pt}
  \item \texttt{finetune/checkpoints/finetuned\_best.pt}
  \item \texttt{finetune/metrics/finetune\_metrics.json}
\end{itemize}

\section{Stage 4: Retraining Oracle}
If enabled, \texttt{scripts/run\_pipeline.py} dispatches \texttt{scripts/train\_vlm.py --oracle}.

\subsection{What changes in oracle mode}
\begin{itemize}
  \item Training starts from the original saved initialization checkpoint, not the fine-tuned checkpoint.
  \item The train loader is forced to \texttt{retain\_train}.
  \item The source fine-tuning metrics are loaded to reconstruct the optimizer-step and scheduler-step budget.
  \item The code validates that the oracle run matches the original architecture and source config.
\end{itemize}

This is the repository's closest approximation to a deletion oracle. It is important because later evaluation can compare unlearned models against this retrained baseline.

\section{Stage 5: Unlearning}
For each configured method, \texttt{scripts/run\_pipeline.py} dispatches \texttt{scripts/run\_unlearning.py}.

\subsection{Wrapper behavior}
\texttt{scripts/run\_unlearning.py} mainly:
\begin{itemize}
  \item resolves the method and all unlearning hyperparameters,
  \item resolves the fine-tuned checkpoint path,
  \item builds an \texttt{UnlearnConfig},
  \item calls \texttt{src/unml/unlearn.py::run\_unlearning}.
\end{itemize}

\subsection{Shared setup inside \texttt{run\_unlearning}}
\begin{enumerate}
  \item Load split metadata and validate the requested method.
  \item For H-TGSD methods, validate the hierarchy request and H-TGSD hyperparameters.
  \item Build CLIP processor, tokenizer, and class text inputs.
  \item Build dataloaders from the same split JSON.
  \item Load the fine-tuned checkpoint as the student model.
  \item Create a frozen deep-copied teacher from that student.
  \item Cache teacher text features, and cache student text features when the text path is frozen.
  \item Create an AdamW optimizer over trainable student parameters.
\end{enumerate}

\subsection{Baseline methods}
For non-H-TGSD methods, the code uses:
\begin{itemize}
  \item \texttt{retain\_only}: supervised CE only on retain data.
  \item \texttt{ga\_kl}: gradient ascent on forget CE plus KL preservation on retain data.
  \item \texttt{counterfactual\_rebind}: push forget samples toward a random wrong class with a margin term, plus retain KL.
  \item \texttt{entropy\_rebind}: maximize entropy on forget samples, plus retain KL.
\end{itemize}

These methods repeatedly draw from:
\begin{itemize}
  \item \texttt{retain\_train}
  \item \texttt{forget}
\end{itemize}

\subsection{H-TGSD-specific setup}
For \texttt{h\_tgsd} and \texttt{h\_tgsd\_no\_sibling\_preservation}, the code first constructs semantic bases:
\begin{enumerate}
  \item Validate request type: only \texttt{superclass} and \texttt{selective\_class} are accepted.
  \item Collect teacher image prototypes per target and sibling class from:
  \begin{itemize}
    \item \texttt{forget}
    \item optionally \texttt{sibling\_retain}
  \end{itemize}
  \item Combine each image prototype with the corresponding CLIP text feature.
  \item Call \texttt{src/unml/disentangle.py::build\_semantic\_bases}.
\end{enumerate}

For selective-class forgetting:
\begin{itemize}
  \item sibling directions span the shared basis,
  \item the target basis is the residual of the target direction after projection off that shared sibling basis.
\end{itemize}

The code also saves basis metadata to:
\[
\texttt{unlearning/metrics/<method>\_basis.pt}
\]

\subsection{H-TGSD optimization loop}
Each unlearning step draws:
\begin{itemize}
  \item a forget batch,
  \item an unrelated retain batch,
  \item and, for full H-TGSD, a sibling retain batch.
\end{itemize}

The helper \texttt{\_student\_teacher\_outputs} computes:
\begin{itemize}
  \item student image features,
  \item student logits,
  \item teacher image features,
  \item teacher logits.
\end{itemize}

Those are passed into \texttt{h\_tgsd\_objective}, whose loss is a weighted sum of:
\begin{itemize}
  \item target-subspace suppression,
  \item forget entropy maximization,
  \item unrelated prediction KL preservation,
  \item unrelated feature preservation,
  \item sibling prediction KL preservation,
  \item sibling feature preservation,
  \item shared-subspace coordinate preservation.
\end{itemize}

The \texttt{h\_tgsd\_no\_sibling\_preservation} ablation still builds the same bases, but disables sibling-preservation behavior in the objective.

\subsection{End of unlearning}
After optimization, the code:
\begin{itemize}
  \item evaluates retain, forget, and test splits,
  \item saves \texttt{unlearn\_<method>.pt},
  \item stores runtime, architecture, metrics, provenance, and optional H-TGSD basis metadata in the checkpoint,
  \item writes \texttt{unlearn\_<method>\_metrics.json}.
\end{itemize}

\section{Stage 6: Evaluation and Attack Comparison}
After all candidate models exist, \texttt{scripts/run\_pipeline.py} dispatches \texttt{scripts/evaluate\_attacks.py}.

\subsection{Internal flow}
\begin{enumerate}
  \item Resolve candidate names and checkpoint paths.
  \item Build an \texttt{AttackConfig}.
  \item Call \texttt{src/unml/attacks.py::run\_attack\_comparison}.
\end{enumerate}

\subsection{What evaluation measures}
The evaluation code compares:
\begin{itemize}
  \item fine-tuned baseline,
  \item retrain oracle if available,
  \item every unlearning method checkpoint.
\end{itemize}

It computes:
\begin{itemize}
  \item target, sibling, unrelated, retain, and all-class accuracy summaries,
  \item forget-train behavior,
  \item confidence-based and loss-based MIA AUC,
  \item optional oracle-relative prediction and embedding distances,
  \item optional H-TGSD subspace-energy metrics when basis metadata is available.
\end{itemize}

\subsection{Why the base checkpoint matters}
The attack suite loads the original base initialization checkpoint and uses it as a reference when forming delta-based membership signals. So evaluation is not only comparing final accuracies; it also compares how confidence or loss changed relative to the base model.

\subsection{Evaluation artifacts}
The pipeline expects:
\begin{itemize}
  \item \texttt{comparison/comparison.csv}
  \item \texttt{comparison/comparison.md}
  \item \texttt{comparison/behavioral\_tradeoff.png}
  \item \texttt{comparison/plot\_manifest.json}
  \item optionally \texttt{comparison/semantic\_subspace.png}
\end{itemize}

\section{Core Shared Infrastructure}
\subsection{Configuration resolution}
\texttt{src/unml/config.py} is the central resolver for:
\begin{itemize}
  \item dataset selection,
  \item request selection,
  \item stage-specific defaults,
  \item output-directory conventions.
\end{itemize}

Important implication: many paths and settings are not hard-coded in the stage scripts. They are derived through config helpers such as \texttt{resolve\_dataset\_and\_split\_path}, \texttt{resolve\_stage\_output\_dir}, and \texttt{resolve\_request\_config}.

\subsection{Data loaders}
\texttt{src/unml/data.py::build\_loaders} creates named loaders for all later stages. These named loaders are the common language of the pipeline:
\begin{itemize}
  \item \texttt{forget}
  \item \texttt{retain\_train}
  \item \texttt{retain\_val}
  \item \texttt{finetune\_train}
  \item \texttt{test\_forget}
  \item \texttt{test\_retain}
  \item \texttt{test\_all}
  \item \texttt{sibling\_retain}
  \item \texttt{unrelated\_retain}
  \item \texttt{test\_sibling}
  \item \texttt{test\_unrelated}
\end{itemize}

\subsection{Evaluation helper usage}
\texttt{src/unml/evaluate.py} provides:
\begin{itemize}
  \item text-feature construction,
  \item standard classification evaluation,
  \item output collection for logits, probabilities, embeddings, and per-sample losses,
  \item grouped summaries by class subsets.
\end{itemize}

\section{Concrete Execution Summary}
In one sentence, the actual pipeline is:

\begin{quote}
Create a hierarchy-aware split JSON, fine-tune a frozen-CLIP adapter model, optionally retrain from the original adapter initialization on retained data only, unlearn from the fine-tuned checkpoint with one or more methods, then compare all resulting checkpoints with utility, forgetting, MIA, oracle-distance, and optional semantic-subspace metrics.
\end{quote}

\section{Practical Gotchas}
\begin{itemize}
  \item H-TGSD depends on named hierarchy requests and sibling metadata in the split file.
  \item H-TGSD requires cacheable text features, which effectively means the text path must remain frozen.
  \item The evaluation stage assumes all declared candidate checkpoints exist before comparison starts.
  \item \texttt{--resume} is receipt-based, not just file-existence-based.
  \item The retraining oracle is only run if enabled in config or explicitly requested.
\end{itemize}

\section{Research Notes Moved from Code}
\subsection{CLIP logit scale during training and unlearning}
\texttt{logit\_scale} is CLIP's learned temperature parameter. It controls how sharp the image--text similarity distribution is before softmax. In pretrained CLIP it is typically close to \(\ln(100)\). Training this scalar gives the adapter another degree of freedom to adjust classification confidence.

For unlearning evaluation this can become a confound. If one method changes \texttt{logit\_scale} more than another, confidence changes may reflect both forgetting and temperature drift. This matters for delta-based membership inference because that attack compares \(\text{confidence}_{after} - \text{confidence}_{base}\). The CIFAR-100 profile currently freezes \texttt{logit\_scale} during unlearning to reduce this ambiguity.

\subsection{Membership inference metrics}
The confidence MIA asks whether true-label confidence separates training members from non-members. AUC \(=0.5\) means no separability under that score, while values far from \(0.5\) indicate stronger membership signal. The delta MIA uses the change from the base adapter checkpoint rather than raw confidence, which is why aligned dataset indices and stable temperature behavior matter.

\subsection{Baseline unlearning objectives}
\texttt{ga\_kl} maximizes the forget-set cross-entropy while preserving retain behavior through KL-to-teacher regularization. It is a direct behavioral forgetting baseline, but it can damage nearby retained classes when the forget and retain classes share features.

\texttt{counterfactual\_rebind} redirects forget samples toward randomly chosen wrong classes and adds a margin so the counterfactual score exceeds the true-class score. This can saturate: once the model confidently predicts the counterfactual class, gradients from the rebind and margin terms shrink and retain KL dominates the objective.

\texttt{entropy\_rebind} instead maximizes entropy over all classes for forget samples. It does not pick a single counterfactual class; it pushes the forget prediction distribution toward uncertainty while retaining KL constrains retained behavior.

\end{document}
